% Judul Capstone
\judul{\textit{SCADA-Based Optimal Load Shedding and Recovery Using Mathematical Optimization Under Generation and Feeder Contingencies}}

% Jenis Dokumen
%% Jenis Dokumen : PERANCANGAN PRODUK DAN SPESIFIKASI

% Kode Dokumen
%% Kode Dokumen : C-251

% Nomor Dokumen (ID Kelompok Capstone)
\NoDok{D\_03}

% Nomor Revisi
\NoRev{01}

% Tanggal Penerbitan Dokumen
%% Otomatis terisi tanggal ketika file LaTeX ini di-compile

% Data Mahasiswa Capstone
%% Format : \MHS{<Nama Lengkap>}{<NIM>}{<Prodi>}{<Alamat Email>}
% Ketua Kelompok
\MHSA{Muhammad Akbar Maulana}{23/514723/TK/56512}
	  {Teknik Elektro}{muhammadakbarmaulana@mail.ugm.ac.id}
% Anggota 1
\MHSB{Muhammad Basel Fawaz Sigit}{23/516761/TK/56809}
	  {Teknik Elektro}{muhammadbaselfawazsigit516761@mail.ugm.ac.id} 
% Anggota 2
\MHSC{Nathaniel Benelisha}{23/517262/TK/56883}
	  {Teknik Elektro}{nathanielbenelisha@mail.ugm.ac.id}
% Anggota 3
\MHSD{Ara Dwi Narendra}{23/522807/TK/57752}
	  {Teknik Elektro}{aradwinarendra@mail.ugm.ac.id}
% Anggota 4
\MHSE{Dimitris G. Manolakis}{20/443100/TK/54200}
	  {Teknik Elektro}{dgmanolakis@mail.ugm.ac.id}
% Un-comment Line di bawah ini apabila tidak ada Anggota 4 :
\MHSE{}{}{}{}

% Dosen Pembimbing
%% Format : {<Nama Lengkap>}{<NIP/NIU>}
\DPA{Dr. Wijaya Yudha Atmaja, S.T., M.Eng.}{111198907202408101}

% Tempat Pelaksanaan.
%% Format : {<Nama Laboratorium> \newline <Nama Departemen> \newline <Nama Fakultas>}
\Tempat{Departemen Teknik Elektro dan Teknologi Informasi \newline Fakultas Teknik}